\section{Prompts used in the experiments}
\label{sec:appendix-prompts}

The full text of the LLM prompts is reproduced here for exact reproduction.
Variables in \code{\{braces\}} are filled at runtime.

\subsection{Compaction prompt (used by all four strategies)}
\label{sec:prompt-compaction}

System message:

\begin{quote}
You are a conversation summarizer. Be concise and precise.
\end{quote}

User message (followed by the conversation segment to compact):

\begin{lstlisting}
Summarize the following conversation segment.

PRESERVE (critical):
- Key facts and decisions made
- File paths and code structures mentioned
- User preferences and requirements stated
- Actions completed (files created/edited, commands run)
- Errors encountered and how they were resolved

DISCARD:
- Full file contents (keep only: filename, line count, key elements)
- Verbose tool outputs (keep only: what was done + result)
- Intermediate reasoning that led nowhere
- Redundant back-and-forth

Format: Concise narrative summary. Bullet points for lists of facts.
Target: ~20% of original length, keeping all actionable information.

---
\end{lstlisting}

\subsection{Q\&A prompt (used in \S\ref{sec:calibration}, \S\ref{sec:singlepass}, \S\ref{sec:strategy} evaluations)}
\label{sec:prompt-qa}

System message:

\begin{quote}
You are a helpful assistant working on a complex software project.
Answer questions precisely from memory.
\end{quote}

User message (sent after the test conversation context):

\begin{lstlisting}
Answer each of the following questions based ONLY on what you know from
our conversation.
Be specific: include exact numbers, names, paths, versions.
If you don't remember or aren't sure, say "I don't recall".

{questions}

Reply with a JSON array of objects, one per question:
[{"id": "LM_0001", "answer": "your answer"}, ...]

IMPORTANT: Return ONLY the JSON array, no other text.
\end{lstlisting}

\subsection{Strict judge prompt (used to score \S\ref{sec:strategy} results)}
\label{sec:prompt-judge}

System message:

\begin{quote}
You are an objective evaluator. Answer ONLY with valid JSON.
\end{quote}

User message:

\begin{lstlisting}
Evaluate each LLM answer against the expected answer using STRICT criteria.

{entries}

For each entry, decide TWO booleans:

1. "recalled": Is the SPECIFIC EXPECTED FACT actually present in the LLM's
   answer?
   - TRUE only if the LLM's answer contains the substantive expected
     information (the actual fact, value, name, or concept asked about).
   - FALSE if any of the following:
     * The answer says "I don't recall" / "I don't know" / "I'm not sure"
     * The answer denies the topic was discussed ("I don't recall you
       mentioning X")
     * The answer mentions a RELATED topic but not the actual expected fact
     * The answer is vague/generic and could match many possible expected
       values
     * The answer contradicts the expected fact

2. "accurate": Is the answer factually equivalent to the expected answer?
   - TRUE if the answer contains the expected value (semantic match,
     exact numbers, etc.)
   - FALSE if values differ, dates differ, or the answer is wrong on the
     substance.
   - Note: accurate=TRUE implies recalled=TRUE.

Examples (STRICT):
- expected "17 days", answer "10 days" -> recalled=TRUE, accurate=FALSE
- expected "turbinado sugar for cookies", answer "I don't recall extras
  for cookies but cookies were made" -> recalled=FALSE
- expected "Miami trip in June", answer "I don't recall any Miami trip"
  -> recalled=FALSE
- expected "Plex Media Server", answer "Plex" -> recalled=TRUE, accurate=TRUE

Reply with ONLY a JSON array:
[{"id": "LM_0001", "recalled": true/false, "accurate": true/false,
  "notes": "brief reason"}, ...]
\end{lstlisting}

A ``lenient'' version of the judge prompt (without the explicit FALSE
clauses for ``I don't recall'' etc.) was used in early \S\ref{sec:strategy} runs and was
later replaced with the strict version above. See \S\ref{sec:methodological} for the impact of
this change. Both versions are in the repository at
\code{benchmark\_iterative\_v6.py} (\code{BATCH\_JUDGE\_PROMPT} and
\code{BATCH\_JUDGE\_PROMPT\_STRICT}).
